\documentclass[aspectratio=169,10pt]{beamer}
\usetheme[numbering=fraction,progressbar=frametitle]{metropolis}
\usepackage{fontspec}
\setsansfont{Latin Modern Sans}
\usepackage{amsmath,amssymb,mathtools,graphicx,booktabs}
\usepackage{polyglossia}
\setdefaultlanguage{french}
\setbeamertemplate{navigation symbols}{}
\setbeamerfont{frametitle}{size=\large}
\setbeamerfont{normal text}{size=\normalsize}
\setlength{\parskip}{.45em}
\hypersetup{pdfauthor={},pdfsubject={Lecture slides},pdfkeywords={},pdftitle={Variations et extrémums d'une fonction}}
\providecommand{\N}{\mathbb{N}}
\providecommand{\Z}{\mathbb{Z}}
\providecommand{\D}{\mathbb{D}}
\providecommand{\Q}{\mathbb{Q}}
\providecommand{\R}{\mathbb{R}}
\providecommand{\C}{\mathbb{C}}
\providecommand{\vect}[1]{\overrightarrow{#1}}
\providecommand{\vi}{\vec{\imath}}
\providecommand{\vj}{\vec{\jmath}}
\providecommand{\vk}{\vec{k}}
\providecommand{\norm}[1]{\left\lVert #1\right\rVert}
\providecommand{\vc}[2]{\begin{pmatrix}#1\\#2\end{pmatrix}}
\providecommand{\deter}[2]{\det\left(#1,#2\right)}
\providecommand{\abs}[1]{\left\lvert #1\right\rvert}
\providecommand{\dd}{\,\mathrm{d}}
\providecommand{\e}{\mathrm{e}}
\providecommand{\Card}{\operatorname{Card}}
\providecommand{\E}{\operatorname{E}}
\providecommand{\V}{\operatorname{V}}
\providecommand{\degre}{^{\circ}}
\providecommand{\un}[1]{\,\mathrm{#1}}
\providecommand{\sci}[2]{#1\times10^{#2}}
\providecommand{\pgcd}{\operatorname{PGCD}}
\providecommand{\ppcm}{\operatorname{PPCM}}
\title{Variations et extrémums d'une fonction}
\subtitle{Cours et exemples}
\date{}
\author{}
\begin{document}
\begin{frame}[plain]
\vfill
{\usebeamerfont{title}\inserttitle\par}
\medskip
{\usebeamerfont{subtitle}\insertsubtitle\par}
\medskip
{\usebeamercolor[fg]{progress bar}\rule{\textwidth}{.5pt}}
\vfill
\end{frame}
\begin{frame}{Croissance, décroissance et monotonie}
Sur un intervalle \(I\), \(f\) est croissante si pour tous \(a,b\in I\), \(a<b\Rightarrow f(a)\leq f(b)\) ; décroissante si \(a<b\Rightarrow f(a)\geq f(b)\). Les variations sont strictes lorsque les inégalités entre images sont strictes. Une fonction constante a la même valeur partout ; elle est à la fois croissante et décroissante au sens large.
Une fonction monotone sur \(I\) y est croissante ou décroissante.
\end{frame}
\begin{frame}{Croissance, décroissance et monotonie}
Indiquer le domaine, les intervalles de monotonie, les valeurs aux bornes lorsqu'elles sont définies, puis les flèches. Un tableau de signes indique si \(f(x)\) est positif ou négatif ; un tableau de variations indique comment ses valeurs évoluent.
\end{frame}
\begin{frame}{Croissance, décroissance et monotonie}
Sur \([-3;4]\), \(f(x)=(x-1)^2-2\). À l'aide du carré, en prolongeant la courbe étudiée en séquence 6, dresser son tableau de variations. Donner deux nombres \(a<b\) tels que \(f(a)>f(b)\).
\end{frame}
\begin{frame}{Correction : Croissance, décroissance et monotonie}
La fonction décroît strictement sur \([-3;1]\), puis croît strictement sur \([1;4]\). \(f(-3)=14\), \(f(1)=-2\), \(f(4)=7\).
\[\begin{array}{c|ccccc}x&-3&&1&&4\\\hline f(x)&14&\searrow&-2&\nearrow&7\end{array}\]
Par exemple \(a=0<b=1\) et \(f(0)=-1>-2=f(1)\).
\end{frame}
\begin{frame}{Croissance, décroissance et monotonie}
Une fonction peut-elle être négative et croissante ? Donner un exemple sur un intervalle.
\end{frame}
\begin{frame}{Correction : Croissance, décroissance et monotonie}
Oui : \(f(x)=x-5\) sur \([0;2]\) est strictement croissante et reste entre \(-5\) et \(-3\), donc strictement négative.
\end{frame}
\begin{frame}{Maximum et minimum sur un intervalle}
\(f\) admet un maximum en \(a\in I\) si, pour tout \(x\in I\), \(f(x)\leq f(a)\). La valeur du maximum est \(f(a)\).
\(f\) admet un minimum en \(b\in I\) si, pour tout \(x\in I\), \(f(x)\geq f(b)\).
Un extrémum doit être atteint. Distinguer sa valeur de l'abscisse où il est atteint.
\end{frame}
\begin{frame}{Maximum et minimum sur un intervalle}
Reprendre \(f(x)=(x-1)^2-2\) sur \([-3;4]\). Donner minimum et maximum, puis leurs lieux d'atteinte. Refaire sur \([0;4]\).
\end{frame}
\begin{frame}{Correction : Maximum et minimum sur un intervalle}
Sur \([-3;4]\), minimum \(-2\) en \(x=1\), maximum \(14\) en \(x=-3\). Sur \([0;4]\), le minimum reste \(-2\) en \(1\), mais le maximum vaut \(7\) en \(4\). Le maximum dépend du domaine.
\end{frame}
\begin{frame}{Maximum et minimum sur un intervalle}
La fonction \(f(x)=x\) possède-t-elle un maximum sur \([0;1[\) ? Un minimum ?
\end{frame}
\begin{frame}{Correction : Maximum et minimum sur un intervalle}
Elle a pour minimum \(0\), atteint en \(0\). Elle n'a pas de maximum : pour tout \(x\in[0;1[\), \((x+1)/2\) appartient encore à \([0;1[\) et possède une image plus grande. Le nombre \(1\) est une borne supérieure, mais n'est pas atteint.
\end{frame}
\begin{frame}{Maximum et minimum sur un intervalle}
Un tableau de variations ne fournit pas toutes les images intermédiaires. Il permet de comparer certaines images ; si les abscisses se situent sur des branches différentes, la comparaison peut être indéterminée.
\end{frame}
\begin{frame}{Maximum et minimum sur un intervalle}
Le seul fait qu'une fonction décroisse puis croisse permet-il de comparer deux images situées de part et d'autre du minimum ?
\end{frame}
\begin{frame}{Correction : Maximum et minimum sur un intervalle}
Non. Il faut des informations supplémentaires. Pour \(x\mapsto x^2\), les images de \(-2\) et \(1\) valent \(4\) et \(1\), tandis que celles de \(-1\) et \(2\) valent \(1\) et \(4\).
\end{frame}
\begin{frame}{Variations des fonctions affines}
Pour \(f(x)=mx+p\) et \(a\ne b\),
\[\dfrac{f(b)-f(a)}{b-a}=m.\]
Si \(m>0\), \(f\) est strictement croissante ; si \(m<0\), strictement décroissante ; si \(m=0\), constante.
\end{frame}
\begin{frame}{Variations des fonctions affines}
Pour \(a<b\), \(f(b)-f(a)=m(b-a)\). Comme \(b-a>0\), la différence a le signe de \(m\). Elle est positive, négative ou nulle selon les trois cas, ce qui donne les variations annoncées. La division par \(b-a\) donne le taux d'accroissement constant.
\end{frame}
\begin{frame}{Variations des fonctions affines}
Déterminer les variations de \(f(x)=3x+2\), \(g(x)=7-6x\), \(h(x)=-x\). Donner le signe de \(2x-6\) et expliquer pourquoi le signe de cette fonction ne reste pas constant alors qu'elle est croissante.
\end{frame}
\begin{frame}{Correction : Variations des fonctions affines}
\(f\) est strictement croissante ; \(g,h\) strictement décroissantes. \(2x-6\) est négatif pour \(x<3\), nul en \(3\), positif pour \(x>3\). Une fonction croissante peut franchir \(0\) ; son sens de variation ne détermine pas un signe unique.
\end{frame}
\begin{frame}{Variations des fonctions affines}
Une fonction affine \(u\) vérifie \(u(2)=5\), \(u(6)=13\). Déterminer \(u\), puis son minimum et son maximum sur \([-1;4]\).
\end{frame}
\begin{frame}{Correction : Variations des fonctions affines}
\(m=(13-5)/(6-2)=2\), \(p=1\), donc \(u(x)=2x+1\). Elle est croissante : minimum \(u(-1)=-1\), maximum \(u(4)=9\).
\end{frame}
\begin{frame}{Variations du carré et de la valeur absolue}
\(x\mapsto x^2\) et \(x\mapsto|x|\) sont strictement décroissantes sur \(]-\infty;0]\) et strictement croissantes sur \([0;+\infty[\). Leur minimum sur \(\R\) vaut \(0\), atteint uniquement en \(0\).
\end{frame}
\begin{frame}{Variations du carré et de la valeur absolue}
Pour \(a<b\), \(b^2-a^2=(b-a)(a+b)\). Si \(0\leq a<b\), les deux facteurs sont strictement positifs : \(a^2<b^2\). Si \(a<b\leq0\), le premier facteur est positif et le second négatif : \(a^2>b^2\).
Pour la valeur absolue, sur \([0;+\infty[\), \(|x|=x\) est croissante ; sur \(]-\infty;0]\), \(|x|=-x\) est décroissante. La positivité et l'annulation en \(0\) donnent le minimum.
\end{frame}
\begin{frame}{Variations du carré et de la valeur absolue}
Comparer sans calculatrice : \((-3{,}2)^2\) et \((-3{,}1)^2\) ; \(1{,}2^2\) et \(1{,}3^2\) ; \(|-5|\) et \(|-2|\).
\end{frame}
\begin{frame}{Correction : Variations du carré et de la valeur absolue}
Sur les négatifs, le carré renverse l'ordre : \((-3{,}2)^2>(-3{,}1)^2\). Sur les positifs, \(1{,}2^2<1{,}3^2\). Enfin \(|-5|>|-2|\), soit \(5>2\).
\end{frame}
\begin{frame}{Variations du carré et de la valeur absolue}
Encadrer \(x^2\) si \(-3\leq x\leq2\). Pourquoi ne peut-on pas simplement élever les deux bornes au carré ?
\end{frame}
\begin{frame}{Correction : Variations du carré et de la valeur absolue}
Le minimum est \(0\) et le maximum \(9\) : \(0\leq x^2\leq9\). Le carré n'est pas monotone sur \([-3;2]\). Élever les seules bornes donnerait à tort \(9\leq x^2\leq4\).
\end{frame}
\begin{frame}{Variations de la fonction inverse}
La fonction inverse est strictement décroissante sur \(]-\infty;0[\) et sur \(]0;+\infty[\), séparément. Elle n'est pas décroissante sur tout \(\R^*\).
\end{frame}
\begin{frame}{Variations de la fonction inverse}
Pour \(a<b\) de même signe, \(ab>0\) et
\[\dfrac1b-\dfrac1a=\dfrac{a-b}{ab}<0.\]
Donc \(1/b<1/a\). En revanche, \(-1<1\) et \(1/(-1)=-1<1=1/1\) : le critère de décroissance échoue sur tout \(\R^*\).
\end{frame}
\begin{frame}{Variations de la fonction inverse}
Comparer \(1/(-4)\) et \(1/(-2)\), puis \(1/3\) et \(1/5\). Encadrer \(1/x\) si \(2\leq x\leq5\), puis si \(-5\leq x\leq-2\).
\end{frame}
\begin{frame}{Correction : Variations de la fonction inverse}
\(-1/4>-1/2\) ; \(1/3>1/5\). Pour \(x\in[2;5]\), \(1/5\leq1/x\leq1/2\). Pour \(x\in[-5;-2]\), \(-1/2\leq1/x\leq-1/5\).
\end{frame}
\begin{frame}{Variations de la fonction inverse}
Pourquoi l'intervalle \([-1;1]\) ne convient-il pas pour appliquer directement une règle de variation de la fonction inverse ?
\end{frame}
\begin{frame}{Correction : Variations de la fonction inverse}
La fonction n'est pas définie en \(0\), qui appartient à l'intervalle. Il faut travailler sur les deux parties du domaine ; de plus, la fonction n'a pas le même comportement de comparaison lorsqu'on traverse \(0\).
\end{frame}
\begin{frame}{Position relative d'un nombre et de son carré}
Pour \(0\leq x\leq1\), \(x^2\leq x\) ; pour \(x\geq1\), \(x^2\geq x\). L'égalité a lieu en \(0\) et \(1\).
\end{frame}
\begin{frame}{Position relative d'un nombre et de son carré}
\(x^2-x=x(x-1)\). Sur \([0;1]\), le premier facteur est positif ou nul, le second négatif ou nul : leur produit est négatif ou nul. Sur \([1;+\infty[\), les deux sont positifs ou nuls. Le produit est nul si \(x=0\) ou \(x=1\).
\end{frame}
\begin{frame}{Position relative d'un nombre et de son carré}
La fonction racine carrée est strictement croissante sur \([0;+\infty[\). La fonction cube est strictement croissante sur \(\R\).
\end{frame}
\begin{frame}{Position relative d'un nombre et de son carré}
Expliquer la croissance de la racine carrée à partir de celle du carré sur les positifs. Pour le cube, utiliser \(b^3-a^3=(b-a)(a^2+ab+b^2)\).
\end{frame}
\begin{frame}{Correction : Position relative d'un nombre et de son carré}
Si \(0\leq a<b\) mais \(\sqrt a\geq\sqrt b\), la croissance du carré sur les positifs donnerait \(a\geq b\), contradiction. Pour \(a<b\), \(b-a>0\) et \(a^2+ab+b^2=(a+b/2)^2+3b^2/4>0\) (l'égalité exigerait \(a=b=0\)). Donc \(b^3-a^3>0\).
\end{frame}
\begin{frame}{Position relative d'un nombre et de son carré}
Comparer \(\sqrt{5}\) et \(\sqrt6\), puis \((-2)^3\) et \((-1)^3\).
\end{frame}
\begin{frame}{Correction : Position relative d'un nombre et de son carré}
\(\sqrt5<\sqrt6\) et \(-8<-1\), en accord avec les deux croissances.
\end{frame}
\begin{frame}{Optimiser : conjecturer et démontrer}
Un rectangle de périmètre \(20\) m a pour aire \(A(x)=x(10-x)\) avec \(x\in]0;10[\).

\begin{enumerate}
\setcounter{enumi}{0}\item Conjecturer l'aire maximale par un tableau de valeurs.
\item Montrer \(A(x)=25-(x-5)^2\). En déduire le maximum exact et les dimensions correspondantes.

\end{enumerate}
\end{frame}
\begin{frame}{Correction : Optimiser : conjecturer et démontrer}
Pour \(x=1,2,3,4,5,6,7,8,9\), les aires sont \(9,16,21,24,25,24,21,16,9\) m\(^2\). En développant, \(25-(x-5)^2=10x-x^2=A(x)\). Comme \((x-5)^2\geq0\), \(A(x)\leq25\), avec égalité seulement en \(x=5\). Le rectangle optimal est un carré de côté \(5\) m.
\end{frame}
\begin{frame}[fragile]{Python}
\small
\begin{verbatim}
def aire(x):
    return x*(10-x)

def maximum_grille(N):
    # N entier >= 2 ; points strictement entre 0 et 10
    meilleur_x = 10/N
    meilleur_y = aire(meilleur_x)
    for k in range(2, N):
        x = 10*k/N
        if aire(x) > meilleur_y:
            meilleur_x, meilleur_y = x, aire(x)
    return meilleur_x, meilleur_y
\end{verbatim}
\end{frame}
\begin{frame}{Optimiser : conjecturer et démontrer}
Que renvoie le programme pour \(N=10\) ? Un balayage numérique prouve-t-il toujours le maximum d'une fonction sur un intervalle ?
\end{frame}
\begin{frame}{Correction : Optimiser : conjecturer et démontrer}
Pour \(N=10\), \((5;25)\). Le programme trouve le maximum sur les seuls points échantillonnés. Il peut manquer un maximum entre deux points : la preuve algébrique donne ici le résultat exact. Une dichotomie ordinaire de recherche d'un zéro ne s'applique pas automatiquement à la recherche d'un extrémum.
\end{frame}
\end{document}
